\section{K-027: an integral exceptional-locus reduction for the exact dyadic pole}

This reduction uses the integral Koszul class of P-F23 before any
normal Euler class is inverted.  Its conclusion is conditional: it
does not prove that the exceptional locus defined below is empty or
that its contribution fails to cancel the chain.

Keep the notation of K-026, with $n\ge4$ and $M=2^{n-1}$, and put
$Y=X_n^{\mu_M}$.  Let $U\subset X_n$ be the open quotient locus of
effectively free closed tensor orbits.  By P-F60, the chain component
$Y_n$ lies in $U$ away from the origin.  Call an irreducible component
$Y_i\ne Y_n$ of $Y_{\rm red}$ \emph{good} if it meets $U$ and
\[
              Y_i\cap X_n^{\mu_{3M}}\subseteq\{0\}.
                                                        \tag{K-027a}
\]
Define the closed, grading- and $S_3$-stable exceptional set
\[
 Z=(Y_{\rm red}\setminus U)\ \cup\!
       \bigcup_{Y_i\ \text{not good},\;Y_i\ne Y_n}Y_i.
                                                        \tag{K-027b}
\]
The first term includes nonfree boundaries of generically free
components.  Thus no hypothesis about such boundaries is suppressed.
The origin belongs to $Z$.

Fix a primitive $M$th root $\xi$ and a prime $\mathfrak p$ above
three in $\mathbb Z[\xi]$.  Write
\[
 \mathcal O=\mathbb Z[\xi]_{\mathfrak p},\qquad
 K=\mathbb Q(\xi),\qquad u=1-t^M.
\]
Laurent expansions below use the branch $t=\xi$ at $u=0$.
Coherent extensions may involve integer grading shifts; their Hilbert
series are bounded-below Laurent series in $t$.  Multiplication by
a power of $t$ reduces any assertion below to a power series and
does not change a pole or integrality statement at $\xi\ne0$.

\begin{proposition}[K-027]
Choose a finite $S_3$-stable homogeneous generating space for
$R_n=\mathbb C[V_n]^{SL_n^3}$, and list the degrees of its generators
not divisible by $M$ as $d_1,\ldots,d_b$.  Set
\[
                         D_M(t)=\prod_{j=1}^b(1-t^{d_j}).
\]
There is an integral virtual graded coherent module class
$\beta$ supported on $Z$, with Hilbert series $H_\beta(t)$,
such that
\[
 H_n^{\rm ten}(t)
       -\frac{H_\beta(t)}{D_M(t)}
       -\frac{C_n(\xi)}{u}
             \ \in\ \mathcal O[[u]]+3\mathcal O((u)).
                                                        \tag{K-027c}
\]
Here $C_n(\xi)$ is the chain's actual normalized local coefficient
from P-F60/P-F23; it is a unit of $\mathcal O$.

In particular, if $H_n^{\rm ten}$ is regular at $\xi$, the principal
part of $H_\beta/D_M$ is integral over $\mathcal O$, its coefficient
of $u^{-1}$ is congruent to $-C_n(\xi)$ modulo three, and every
higher negative Laurent coefficient is divisible by three.
Thus cancellation of the dyadic pole requires a nonzero exceptional
contribution.  If $Z$ is supported at the origin, the tensor Hilbert
series has an actual pole at $\xi$.
\end{proposition}

\begin{proof}
We first record two integral facts, with their precise scope.

\emph{No-collision integrality for supported modules.}
Let $Q$ be a finitely generated graded module over a positively
graded affine algebra, and suppose its support has no nonzero point
fixed by $\mu_{M3^a}$ for any $a\ge1$.  Its Hilbert series is rational
with integer coefficients and root-of-unity poles.  None of the
cyclotomic factors $\Phi_{M3^a}$ occurs in its reduced denominator.
Indeed, at a root of order $M3^a$, take the Koszul complex on algebra
generators whose degrees are not divisible by that order.  The
homology is supported on the corresponding fixed locus inside
$\operatorname{Supp}Q$, hence at the origin, and has finite length.
The alternating Hilbert identity has a Laurent-polynomial right side and a
product of factors nonzero at the root on the left.  It proves
regularity there.  This argument includes nilpotent support and does
not require the algebra or module to be reduced.

Normalize the actual reduced denominator as a product of cyclotomic
polynomials with constant term one.  Multiplication by the bounded-below
integer Laurent series makes the resulting Laurent-polynomial numerator
integral.
Modulo a prime above three, a cyclotomic factor can vanish at the
reduction of $\xi$ only if its index is $M3^a$: this follows from
\[
 \overline{\Phi_{e3^a}}
   =\overline{\Phi_e}^{\varphi(3^a)}\quad(a\ge1,\;3\nmid e),
\]
and the exact order $e$ of the roots of $\overline{\Phi_e}$.
After the excluded factors are removed, all denominator factors
other than $\Phi_M$ are units at $\xi$ in $\mathcal O$.
Moreover $(1-t^M)/\Phi_M(t)$ is a unit there.  Consequently every
Laurent coefficient of $H_Q$ at $u=0$ lies in $\mathcal O$.
One may first expand in $t-\xi$; the formal change to $u$ preserves
integrality because its derivative $-M\xi^{M-1}$ is a unit.
The same conclusions hold for integral virtual module classes.

For a good $Y_i$, condition (K-027a) excludes every higher
$\mu_{M3^a}$ as well, since its fixed locus is contained in the
$\mu_{3M}$ fixed locus.  Thus every graded coherent module supported
on $Y_i$ has the asserted integral Laurent expansion.  The same
holds for the chain closure: all its nonzero points have grading
stabilizer exactly $\mu_M$ by P-F60.

\emph{Integral coherent localization.}
For a grading-stable closed subset $Z$ of a noetherian graded affine
scheme $Y$, the degree-zero coherent Grothendieck localization
sequence says that restriction to $Y\setminus Z$ is surjective on
graded coherent classes and that its kernel is generated by classes
supported on $Z$.  Only this elementary form is needed here.
For clarity, coherent extensions exist by covering the open set with
finitely many homogeneous principal opens and clearing the finitely
many denominators in generators and relations.  Morphisms and the
finitely many exact sequences defining a relation can be extended
the same way; the kernels and cokernels by which exactness fails are
supported on the complement.  This proves both assertions for the
coherent Grothendieck group.  The construction preserves the grading.
It also permits an extension from an open subset of a single closed
component to be chosen with support on that component.

Take a homogeneous polynomial algebra on the chosen generating
space, mapping equivariantly onto $R_n$.  The Koszul complex of the
$b$ moving generators is $S_3$-equivariant.  Its homology modules
$Q_j$ are finitely generated, graded, and supported on $Y$, since
every moving generator acts null-homotopically.  The integral class
\[
                          \alpha=\sum_j(-1)^j[Q_j]
\]
is $S_3$-invariant, and the exact Hilbert identity is
\[
                       D_M(t)H_n^{\rm ten}(t)=H_\alpha(t).
                                                        \tag{K-027d}
\]

On $U$, the quotient is smooth.  The fixed scheme of a finite-order
automorphism in characteristic zero is smooth there, so different
components of $Y$ cannot meet on $U$.  Hence $Y\setminus Z$ is a
disjoint union of the chain minus its origin and open subsets of
good components.  By K-026, the latter are permuted in triples by
$\tau=(123)$.  Choose one representative from each triple.
Coherent localization supplies an integral extension $\gamma$ of
the restriction of $\alpha$ to the chain orbit, supported on $Y_n$,
and an integral extension $\eta_i$ for each chosen good component,
supported on $Y_i$.  Since $\alpha$ is invariant, the classes
$\eta_i,\tau\eta_i,\tau^2\eta_i$ extend its three restrictions.
The difference is supported on $Z$.  In the coherent Grothendieck
group this gives
\[
 \alpha=\gamma+
       \sum_i(\eta_i+\tau\eta_i+\tau^2\eta_i)+\beta,
 \qquad \operatorname{Supp}\beta\subseteq Z.
                                                        \tag{K-027e}
\]
The three translates have the same Hilbert series because $\tau$
preserves every degree.  Thus
\[
 H_n^{\rm ten}
       =\frac{H_\gamma}{D_M}
          +3\sum_i\frac{H_{\eta_i}}{D_M}
          +\frac{H_\beta}{D_M}.
                                                        \tag{K-027f}
\]

Every $1-\xi^{d_j}$ is a unit in $\mathcal O$, because it is a
nonzero dyadic Euler factor.  For example, if $\eta$ has nontrivial
two-power order $q$, then
$(1-\eta)(1+\eta+\cdots+\eta^{q/2-1})=2$.
Therefore $D_M(t)$ is invertible in $\mathcal O[[u]]$.  By the
no-collision argument, each good-component term
$H_{\eta_i}/D_M$ belongs to $\mathcal O((u))$.

The chain term is supported on a curve, so it has at most a simple
pole.  Its leading coefficient is determined by the restriction of
the Koszul class to the nonzero chain orbit, independently of an
extension over the origin.  P-F23 and the full smooth chain
calculation P-F60 give
\[
                 \frac{H_\gamma}{D_M}
                      =\frac{C_n(\xi)}{u}+R(u),
                 \qquad R(u)\in\mathcal O[[u]].
\]
The integrality of the regular remainder follows from the same
no-collision argument for the chain closure.  Its density is one,
and its moving tangent Euler quotient is a quotient of nonzero
dyadic factors.  Hence $C_n(\xi)\in\mathcal O^*$.
Substitution in (K-027f) proves (K-027c).

If $H_n^{\rm ten}$ is regular, comparison of the negative Laurent
coefficients gives exactly the asserted restrictions on the
exceptional term.  If $Z$ is supported at the origin, every class
supported there has finite-length Hilbert polynomial, so
$H_\beta/D_M$ is regular at $\xi$.  Its simple Laurent coefficient
is then zero, contradicting its required congruence to the unit
$-C_n(\xi)$ if $H_n^{\rm ten}$ were regular.  This proves the
conditional final assertion.
\end{proof}

\paragraph{Scope of the reduction.}
The exceptional class is not claimed to be canonical: choices of
coherent extensions can move classes supported on $Z$ between the
displayed terms.  The support restriction and the necessary
principal-part congruence are valid for the resulting integral
decomposition.  The argument treats all lower Laurent coefficients,
singular boundaries, and nilpotent fixed structure through the
Koszul module class.  It does not replace the missing analysis of
$Z$ by an integrality assumption on leading Hilbert densities.
